\section{Methodology}
\label{sec:methodology}

Our approach follows the standard event-study tradition
\citep{mackinlay1997event} adapted to a high-frequency intraday
setting. The core idea is to compare price behaviour in a short
window around each news event to a matched baseline drawn from
non-event periods of the same asset. We pay particular attention to
three design decisions that are often glossed over in shorter-horizon
studies: event-time alignment with respect to the minute bar that
contains the event timestamp, the appropriate proxy for USD strength
when only EUR/USD is observed, and the dependency structure induced
by news clusters.

\subsection{Event-time alignment}
\label{ssec:event-time-alignment}

Each event arrives with a UTC timestamp recorded by Discord at
message publication. To avoid contamination of the post-event
window by the fractional bar that contains the event itself, we use
ceiling/floor alignment on the one-minute bar grid. For an event with
timestamp $t_i$:
\begin{equation}
\text{post-event start}(t_i) = \lceil t_i \rceil_{1\text{m}},
\qquad
\text{pre-event end}(t_i) = \lfloor t_i \rfloor_{1\text{m}} - 1\,\text{minute}.
\label{eq:alignment}
\end{equation}
A forward window of $W$ minutes uses the bars
$\lceil t_i \rceil_{1\text{m}}, \ldots, \lceil t_i \rceil_{1\text{m}} + (W{-}1)\,\text{min}$,
all of which are completely post-event. A pre-event window of $W$
minutes ends at $\lfloor t_i \rfloor_{1\text{m}} - 1\,\text{min}$.
Events whose window cannot be completely populated (missing bars,
weekend boundary, market closure) are flagged as
\texttt{is\_in\_closed\_period} and diverted to a separate gap-test
pipeline rather than being silently dropped.

This alignment is more conservative than the common practice of
using $\lfloor t_i \rfloor_{1\text{m}}$ as the base bar, which mixes
the bar containing the event into both the pre-event and post-event
windows and biases short-window estimates toward zero.

\subsection{Return definitions and the USD proxy convention}
\label{ssec:returns-and-usd-proxy}

For each event $i$, asset $a \in \{\text{EUR/USD}, \text{NDX}\}$,
and window length $W \in \{1, 5, 15, 60, 240\}$ minutes, we define
the close-to-close percentage return
\begin{equation}
r^{\text{price}}_{i,a,W} \;=\; \frac{P^{\text{close}}_{\text{end}} - P^{\text{open}}_{\text{start}}}{P^{\text{open}}_{\text{start}}} \times 100,
\label{eq:price-return}
\end{equation}
together with two intra-window magnitude measures,
\begin{align}
\text{range}_{i,a,W} &= \frac{\max_{t \in W} H_t \;-\; \min_{t \in W} L_t}{P^{\text{open}}_{\text{start}}} \times 100,
\label{eq:range} \\
\text{maxabs}_{i,a,W} &= \max\bigl(|\text{max\_up}|, |\text{max\_down}|\bigr),
\label{eq:maxabs}
\end{align}
where $H_t, L_t$ are the high and low of bar $t$ within the window.
The range and maximum-absolute-move statistics capture intra-window
volatility that close-to-close returns can miss, particularly on
news where the price overshoots and partially reverses.

Because the sentiment labels are constructed with respect to USD
strength rather than EUR weakness, the EUR/USD price return must be
sign-flipped before being compared against \texttt{sentiment\_usd}:
\begin{equation}
\rtarget_{i,\text{EUR/USD},W} = -\,r^{\text{price}}_{i,\text{EUR/USD},W},
\qquad
\rtarget_{i,\text{NDX},W} = r^{\text{price}}_{i,\text{NDX},W}.
\label{eq:usd-proxy}
\end{equation}
This convention is preserved across H2, H3, H4, and H14. Treating
EUR/USD as a direct USD target produces the wrong sign on every
directional test and is, in our view, one of the most easily missed
sources of error in this kind of analysis.

\subsection{Event clustering}
\label{ssec:clustering}

News events do not arrive independently. Central bank decisions are
typically followed by press conferences; geopolitical escalations
provoke a cascade of related headlines. A naive treatment that
considers each headline as an independent observation will overstate
the number of effective trials and inflate $t$-statistics through
within-cluster correlation. We address this in two complementary
ways.

First, we assign cluster identities by gap rule: two consecutive
events belong to the same cluster if and only if they are separated
by no more than 15 minutes.\footnote{We treat the 15-minute
threshold as a design parameter. The Online Appendix reproduces the
H1 magnitude result with gap thresholds in $\{5, 10, 15, 30\}$
minutes; the central conclusions are insensitive to the choice
within this range.} At the 24{,}490-row event-window panel this
collapses to 14{,}050 cluster-window rows. Non-regression tests
(H1, H2, C1, etc.) use the first event per cluster as the
independent unit, which yields a more conservative effective sample
size than the raw event count.

Second, regression-based tests (H3, H4, C6) employ cluster-robust
standard errors with the cluster identifier as the grouping variable
\citep{white1980heteroskedasticity}, or HC3 standard errors when the
cluster count is insufficient.\footnote{We use HC3 rather than HC0
or HC1 because HC3 is less anti-conservative at the moderate sample
sizes of our category-level tests; the difference between the two
is negligible in the cells with several hundred clusters.} The two
devices together ensure that within-burst correlation does not
contaminate the reported $p$-values.

A separate cluster-level panel
(\texttt{cluster\_event\_study\_windows.csv}) aggregates event-level
features within each cluster: the modal sentiment, the maximum
surprise level, the dominant category, the count of headlines, and
the mean headline length. This permits cluster-level direction
tests (C1) that ask whether the consensus sentiment of the cluster
predicts the realised target direction.

\subsection{Matched baseline}
\label{ssec:matched-baseline}

To contextualise event-window returns, we draw a baseline pool of
non-event windows of the same length for each asset. The pool
excludes all bar-start timestamps falling within $\pm 60$ minutes of
any event, which removes the obvious correlation between the
baseline and the events it is meant to control for. From this pool
we sample 30 matched baseline returns per event, drawing where
possible from the same hour-of-day and day-of-week bucket. When a
bucket contains fewer than ten observations, the sampler falls back
to the hour-of-day bucket, then to the unconditional pool.\footnote{The
choice of 30 draws per event balances Monte Carlo error against
runtime; doubling the draw count changes the central H1 ratios by
less than $0.5\%$ in the cells we re-ran as a check.} This
stratification absorbs known weekly and intraday volatility
seasonality (e.g., the New York open, the Asia-Europe handover)
without removing the event signal.

For each event-window outcome (close-to-close return, absolute
return, range, max-absolute-move) we report
\begin{equation}
z_{i,a,W} \;=\; \frac{x_{i,a,W} - \mu^{\text{base}}_{a,W,h(i),d(i)}}{\sigma^{\text{base}}_{a,W,h(i),d(i)}},
\label{eq:abnormal-z}
\end{equation}
where $h(i)$ and $d(i)$ are the hour and day-of-week of the event
timestamp. Standardising by the matched baseline second moment makes
the EUR/USD and NDX event-window panels directly comparable on a
common dimensionless scale.

\subsection{Hypothesis tests}
\label{ssec:hypothesis-tests}

The pipeline runs fourteen pre-specified hypotheses (H1--H14) and
eight extensions (C1--C8). For reasons of space we report the main
results in Section~\ref{sec:results} only for those tests that
produced non-trivial findings; the full set is reproduced in the
Online Appendix for auditability. Each test is summarised below;
pre-specified $p$-value definitions are documented in
\texttt{event\_study.py} and are not adjusted post hoc.

The main results in Section~\ref{sec:results} draw on eight
hypotheses. \textbf{H1 (volatility)} tests the absolute event-window
return against the matched baseline by a one-sided Mann-Whitney
$U$ test \citep{mann1947test}, with the Welch $t$-test
\citep{welch1947generalization} reported as a robustness check in
the Online Appendix. \textbf{H2 (direction)} tests the hit rate of
the LLM sentiment against the realised target direction by a
one-sided binomial against $p_0 = 0.5$ and reports Wilson 95\%
confidence intervals. \textbf{H3 (sentiment $\times$ prior trend)}
estimates a cluster-robust OLS of the absolute return on sentiment,
the prior 60-minute trend, and their interaction. \textbf{H4
(closed-period gap)} regresses the realised opening-bar gap on
aggregated target sentiment across each closed period (weekend,
holiday) with HC3 robust standard errors. \textbf{H8 (pre-event
drift)} compares the absolute pre-event return to the matched
baseline; we interpret this as evidence on the timing of
information arrival, with explicit caveat that the Discord
timestamp is not the primary-source timestamp. \textbf{H9
(persistence)} measures the sign agreement between the
$+15$-minute and $+4$-hour returns. \textbf{H11 (time-of-day)}
tests hour-of-day and day-of-week effects on the event/baseline
ratio. \textbf{H14 (cross-asset)} reports the correlation between
\eurusd{} and \ndx{} event-window returns both raw and under the
USD-proxy convention.

We extend the battery with seven additional tests. \textbf{C1} is
the cluster-level analogue of H2. \textbf{C2} compares the range and
maximum-absolute-move outcomes to the matched baseline; these
outcomes have the most uniform statistical signal across cells in
the battery and are reported alongside H1. \textbf{C3} reports the
standardised abnormal $z$-scores (eq.~\ref{eq:abnormal-z}) for all
four outcomes. \textbf{C4} runs per-category targeted tests for
central-bank, geopolitical, political, energy, and corporate news,
replacing the redundant omnibus categorical ANOVA originally
proposed as H7. \textbf{C5} assesses stability of the central
effects across the pre- and post-15~January~2026 split; this date
is the knowledge cutoff of the sentiment LLM and the split lets us
address the memorisation concern directly. \textbf{C6} estimates a
multivariate OLS of the standardised abnormal magnitude on category
dummies, surprise level, expected magnitude, confidence, cluster
size, headline length, and pre-event move, with cluster-robust
standard errors. \textbf{C7} reports winsorisation at the 1\% tails
as a robustness check against extreme observations.

A further six tests are reported only in Section~\ref{sec:limitations}
or the Online Appendix because they produced null or exploratory
findings. The LLM auxiliary labels for expected magnitude and
surprise level (originally H5 and H13) are consolidated as a single
Online Appendix subsection. The LLM confidence calibration
(H6, Brier score) is reported in Section~\ref{sec:limitations} as a
null finding. The categorical omnibus ANOVA (H7) is subsumed by C4
and not reported separately. The volume test (H10) is exploratory
because the Dukascopy NDX series is tick volume rather than
consolidated exchange volume. The bull-versus-bear asymmetry test
(H12) is reported as a null finding. Finally, the cross-model
consensus check (C8) on a 200-event subsample comparing
\texttt{deepseek-v4-flash} and \texttt{deepseek-v4-pro} is
exploratory due to subsample size.

\subsection{Identification strategy}
\label{ssec:identification}

We do not claim causal identification of the news effect in the
structural sense of an instrumental-variable or
regression-discontinuity design. The identifying assumption that
underlies the comparison in Sections~\ref{ssec:matched-baseline}
and~\ref{ssec:hypothesis-tests} is the following: conditional on
asset, window length, hour-of-day, and day-of-week, the
distribution of price behaviour in non-event windows provides a
valid counterfactual for the distribution that would have obtained
in the event windows had the news not arrived. This assumption is
defensible in a context where (i)~the intraday volatility cycle is
absorbed by the hour-of-day and day-of-week stratification,
(ii)~there is no systematic mean drift in the matched-baseline
returns across the sample period, and (iii)~the event-window
selection rule (red-dot or \textsc{breaking}, excluding scheduled
macro releases) does not select non-randomly on the unobserved
component of price behaviour.

We assess threats to each leg of this assumption directly:
condition~(i) is tested by reporting hour-of-day effects on the
event/baseline ratio (H11); condition~(ii) is tested by the
pre-cutoff/post-cutoff split (C5); and condition~(iii) is the
hardest to verify, since it depends on properties of the
publisher's filter that we cannot observe directly. The
pre-event drift documented in Section~\ref{ssec:results-timing}
indicates that the public publication time is not the
informational event time, which is a different identification
issue we discuss in
Section~\ref{sec:limitations}.

\subsection{Statistical power}
\label{ssec:power}

The effective sample size after cluster-dedupe is at most $1{,}405$
clusters and varies by hypothesis cell from $616$ to $1{,}289$
clusters because of the cells that are dropped when an asset window
cannot be fully populated. At an effective $n = 1{,}000$ clusters
and the one-sided $\alpha = 0.05$ level, the binomial test of
direction (H2) has $80\%$ power to detect a hit rate of $52.6\%$
against the $50.0\%$ null, or equivalently a true edge of
$2.6$ percentage points. At the more conservative
post-Benjamini-Hochberg $q$-threshold of $0.05$ implied by our
hypothesis family, the minimum detectable edge widens to roughly
$3.0$--$3.5$ percentage points. For the magnitude test (H1), at
$n = 1{,}000$ events versus $30{,}000$ matched-baseline draws, the
one-sided Mann-Whitney $U$ test has $80\%$ power against an
event/baseline mean ratio of approximately $1.10$, well below the
$1.43$ to $1.87$ ratios we report. The magnitude battery is
therefore comfortably overpowered relative to the effect sizes
observed; the direction battery is right-sized for the effect
magnitude consistent with prior literature
\citep{tetlock2007giving,heston2017news} and underpowered to detect
edges smaller than approximately $3$ percentage points.

\subsection{Multiple testing}
\label{ssec:multiple-testing}

The full hypothesis battery generates more than one hundred
individual $p$-values across the main tests, the auxiliary tests,
and the cells of stratified analyses. Reporting raw $p$-values at
conventional thresholds would substantially overstate the level of
evidence. We therefore apply the Benjamini-Hochberg false discovery
rate procedure \citep{benjamini1995controlling} across the union of
all $p$-values produced by the pipeline (all columns prefixed by
\texttt{p\_} in the \texttt{h*\_results.csv} outputs) and report the
resulting $q$-values alongside the raw $p$-values. The primary
conclusions of Section~\ref{sec:results} are drawn exclusively from
$q$-values; raw $p$-values are reported for transparency.

\subsection{Sentiment validation}
\label{ssec:sentiment-validation}

The LLM sentiment labels are the central novel input of the
analysis. To make the LLM choice defensible, we compare it against
two baselines and reserve a held-out validation sample for human
adjudication.

\paragraph{Dictionary baseline.} We re-classify every event using
the Loughran-McDonald finance sentiment dictionary
\citep{loughran2011liability}. The dictionary produces a single
polarity per document, which we threshold at
$|\text{polarity}| \geq 0.1$ to recover the bull/bear/neutral
classes. Because the dictionary has no asset awareness, the
resulting USD and NDX classifications are identical by construction.
We then re-run the H2-equivalent direction test on the LM labels
and compare hit rates side by side with the LLM labels
(Section~\ref{ssec:results-direction}).

\paragraph{Neural baseline.} We additionally classify every event
with FinBERT-tone \citep{yang2020finbert}, a BERT-family model
pre-trained on a 4.9-billion-token corpus of financial
communications. Like the dictionary baseline, FinBERT outputs a
single polarity (positive/negative/neutral) per document and cannot
produce asset-specific sentiment.

\paragraph{External benchmark validation.} The two reference
baselines above test the LLM relative to alternative classifiers on
our specific event sample, but they do not validate the LLM in
absolute terms against a human-annotated gold standard. To address
this, we re-run all three classifiers on the Financial PhraseBank
benchmark \citep{malo2014good}, the de facto standard finance
sentiment evaluation dataset and the same benchmark
\citet{yang2020finbert} use to validate FinBERT. Financial
PhraseBank contains 4{,}840 sentences from financial news annotated
positive/negative/neutral by 16 finance professionals; we use the
\texttt{sentences\_50agree} configuration (sentences for which at
least half of annotators agree on the label). For DeepSeek, we use a
benchmark-adapted prompt that maps to PhraseBank's three-class
scheme so that the production prompt's USD/NDX framing does not
bias the comparison. Section~\ref{ssec:results-direction} reports
accuracy, macro-$F_1$, per-class $F_1$, and Cohen's $\kappa$ for
each classifier against the gold human consensus. This procedure
validates the classifier on an external, peer-reviewed dataset with
known difficulty levels and without requiring human annotation
effort on our event sample.

\paragraph{Cross-model consensus.} A 200-event subsample is
independently classified by \texttt{deepseek-v4-flash} and
\texttt{deepseek-v4-pro} to measure intra-LLM agreement; this is
reported as exploratory due to sample size.

\paragraph{Human validation (future work).} A separate sample of 200
events, stratified across categories and split pre and post the LLM
knowledge cutoff, is reserved for double-coded expert annotation in
a follow-up study. The annotation infrastructure
(\texttt{prepare\_manual\_validation.py},
\texttt{score\_manual\_validation.py}) is in place. Until that work
completes, the external-benchmark validation described above and the
multi-classifier comparison on our sample serve as the primary
quality evidence for the LLM labels.
